\documentclass[a4paper,12pt]{article}
\usepackage[T2A]{fontenc}
\usepackage[utf8]{inputenc}
\usepackage[russian]{babel}
\usepackage{amsmath,amssymb,mathtools}
\usepackage{PTSans}
\renewcommand{\familydefault}{\sfdefault}
\usepackage{icomma}
\usepackage[a4paper,margin=20mm,headheight=27pt,headsep=10mm,footskip=12mm]{geometry}
\usepackage{microtype}
\usepackage{xcolor}
\usepackage{tikz}
\usetikzlibrary{calc,arrows.meta,positioning,shapes.geometric}
\usepackage{booktabs,tabularx}
\usepackage{enumitem}
\usepackage{parskip}
\usepackage{fancyhdr}
\usepackage{setspace}
\usepackage[hidelinks]{hyperref}

\definecolor{accent}{HTML}{2A7775}
\definecolor{darkaccent}{HTML}{1B5553}
\definecolor{textgray}{HTML}{303536}
\definecolor{lightgray}{HTML}{697476}
\definecolor{pale}{HTML}{E9F2F0}
\color{textgray}
\setstretch{1.24}
\setlength{\parindent}{0pt}
\setlength{\parskip}{0.55em}
\setlist[itemize]{leftmargin=1.4em,topsep=0.15em,itemsep=0.15em}
\setlist[enumerate]{leftmargin=1.65em,topsep=0.3em,itemsep=0.48em}
\pagestyle{fancy}
\fancyhf{}
\fancyhead[L]{\small\color{darkaccent}Числовые промежутки}
\fancyhead[R]{\small\color{darkaccent}Кирилл Зиновьев}
\fancyfoot[C]{\small\color{lightgray}\thepage}
\renewcommand{\headrulewidth}{0.4pt}
\renewcommand{\headrule}{\hbox to\headwidth{\color{accent}\leaders\hrule height\headrulewidth\hfill}}
\newcommand{\headingrule}{\vspace{-0.5em}{\color{accent}\hrule height 0.55pt}\vspace{0.5em}}
\newcommand{\key}[1]{\textbf{\color{darkaccent}#1}}
\newcommand{\axes}[1]{\begin{tikzpicture}[x=0.65cm,y=0.65cm,baseline=-0.15cm]
  \draw[->,line width=.5pt,lightgray] (-4.2,0)--(4.35,0) node[right,scale=.78] {$x$};
  #1
\end{tikzpicture}}
\newcommand{\dotfull}[1]{\fill[accent] (#1,0) circle (2.6pt);}
\newcommand{\dotopen}[1]{\fill[white] (#1,0) circle (2.6pt);\draw[accent,line width=1.05pt] (#1,0) circle (2.6pt);}
\tikzset{fcstart/.style={ellipse,draw=darkaccent,thick,align=center,minimum width=36mm,minimum height=10mm,fill=pale},
 fcact/.style={rounded corners=2pt,rectangle,draw=darkaccent,align=center,thick,minimum width=49mm,minimum height=12mm,text width=48mm,inner sep=3.5pt},
 fcdecision/.style={diamond,aspect=2.15,draw=darkaccent,thick,align=center,inner sep=2pt,text width=32mm},
 fcline/.style={-Latex,thick,draw=accent}}

\begin{document}
% --- Титульный лист ---
\thispagestyle{empty}
\begin{center}
\vspace*{26mm}
{\Large\color{darkaccent}\bfseries ЧЕК-ЛИСТ}\\[11mm]
{\huge\bfseries Числовые промежутки}\\[5mm]
{\Large Неравенства, числовая ось и скобки}\\[13mm]
{\large Математика, 7 класс}\\[3mm]
{\large 14 октября 2026 года}\\[19mm]
{\LARGE\bfseries Кирилл Зиновьев}\\[14mm]
\begin{minipage}{0.78\linewidth}
\centering
Как переходить от рисунка на числовой оси к неравенству и записи промежутка; как проверять границы и переводить текстовую задачу на язык уравнения.
\end{minipage}
\vfill
{\large Лёвин Артём Александрович}\\[2mm]
{\normalsize эксклюзивно для Кирилла Зиновьева}\\[13mm]
\begin{tikzpicture}
\draw[accent,line width=1.2pt]
(-6.05,0) .. controls (-4.05,.65) and (-1.9,-.52) .. (0,0)
.. controls (2.2,.67) and (4.3,-.42) .. (6.05,.13);
\end{tikzpicture}
\end{center}
\clearpage

\section*{1. Что означает штриховка на числовой оси}
\headingrule
На оси отмечают \key{все подходящие значения $x$}. Штриховка вправо от точки означает числа больше границы, влево --- числа меньше границы. Граница может \emph{входить} или \emph{не входить} в промежуток.

\begin{center}
\renewcommand{\arraystretch}{1.55}
\begin{tabularx}{\linewidth}{@{}p{63mm}X@{}}
\toprule
\textbf{Рисунок на оси} & \textbf{Неравенство и промежуток}\\
\midrule
\axes{\draw[accent,line width=2.4pt] (-4.05,0)--(1.2,0);\dotfull{1.2}\node[below,font=\scriptsize] at (1.2,-.1) {$5$};}
& $x\le 5$ \quad $x\in(-\infty;5]$\\
\axes{\draw[accent,line width=2.4pt] (-4.05,0)--(1.2,0);\dotopen{1.2}\node[below,font=\scriptsize] at (1.2,-.1) {$5$};}
& $x<5$ \quad $x\in(-\infty;5)$\\
\axes{\draw[accent,line width=2.4pt] (-1.2,0)--(4.1,0);\dotfull{-1.2}\node[below,font=\scriptsize] at (-1.2,-.1) {$5$};}
& $x\ge 5$ \quad $x\in[5;+\infty)$\\
\axes{\draw[accent,line width=2.4pt] (-1.2,0)--(4.1,0);\dotopen{-1.2}\node[below,font=\scriptsize] at (-1.2,-.1) {$5$};}
& $x>5$ \quad $x\in(5;+\infty)$\\
\bottomrule
\end{tabularx}
\end{center}
\key{Как читать:} $x\in[5;+\infty)$ --- «$x$ принадлежит промежутку от $5$ включительно до плюс бесконечности».

\key{Запомните связь трёх обозначений:}
\begin{itemize}
\item Закрашенная точка $\longleftrightarrow$ знаки $\le$, $\ge$ $\longleftrightarrow$ квадратная скобка $[\,]$.
\item Выколотая точка $\longleftrightarrow$ знаки $<$, $>$ $\longleftrightarrow$ круглая скобка $(\,)$.
\item При $-\infty$ и $+\infty$ скобка \textbf{всегда круглая}: бесконечность не является числом, которое можно включить.
\end{itemize}

\key{Важно.} Рисунки схематичны: место числа $5$ на каждом рисунке выбрано для наглядности. Значение границы задаёт подпись под точкой, направление --- штриховка.
\clearpage

\section*{2. Промежуток между двумя числами}
\headingrule
Если штриховка находится \key{между двумя границами}, используйте \key{двойное неравенство}. Слева записывайте меньшую границу, справа --- большую.

\begin{center}
\renewcommand{\arraystretch}{1.68}
\begin{tabularx}{\linewidth}{@{}p{63mm}X@{}}
\toprule
\textbf{Границы на оси} & \textbf{Запись}\\
\midrule
\axes{\draw[accent,line width=2.4pt] (-2.4,0)--(2.5,0);\dotfull{-2.4}\dotfull{2.5}\node[below,font=\scriptsize] at (-2.4,-.1) {$-2$};\node[below,font=\scriptsize] at (2.5,-.1) {$2{,}5$};}
& $-2\le x\le 2{,}5$\\[-1mm]
& $x\in[-2;2{,}5]$\\
\addlinespace[2mm]
\axes{\draw[accent,line width=2.4pt] (-2.4,0)--(2.5,0);\dotopen{-2.4}\dotopen{2.5}\node[below,font=\scriptsize] at (-2.4,-.1) {$-2$};\node[below,font=\scriptsize] at (2.5,-.1) {$2{,}5$};}
& $-2<x<2{,}5$\\[-1mm]
& $x\in(-2;2{,}5)$\\
\addlinespace[2mm]
\axes{\draw[accent,line width=2.4pt] (-2.4,0)--(2.5,0);\dotfull{-2.4}\dotopen{2.5}\node[below,font=\scriptsize] at (-2.4,-.1) {$4$};\node[below,font=\scriptsize] at (2.5,-.1) {$12$};}
& $4\le x<12$\\[-1mm]
& $x\in[4;12)$\\
\addlinespace[2mm]
\axes{\draw[accent,line width=2.4pt] (-2.4,0)--(2.5,0);\dotopen{-2.4}\dotfull{2.5}\node[below,font=\scriptsize] at (-2.4,-.1) {$4$};\node[below,font=\scriptsize] at (2.5,-.1) {$12$};}
& $4<x\le 12$\\[-1mm]
& $x\in(4;12]$\\
\bottomrule
\end{tabularx}
\end{center}

\key{Проверка по границам.} Если $x\in[4;12)$, число $4$ подходит, число $12$ --- нет; при этом число $11{,}99$ подходит. Между двумя разными числами лежит бесконечно много действительных чисел.

\key{Почему в скобках точка с запятой?} В русской записи $2{,}5$ --- одно десятичное число. Точка с запятой помогает однозначно разделять \emph{границы} промежутка.

\section*{3. Мини-чек-лист перед записью ответа}
\headingrule
\begin{itemize}
\item[$\square$] Я вижу, какая часть оси выделена: влево, вправо или между границами.
\item[$\square$] Я выяснил, входит ли каждая конечная граница.
\item[$\square$] Я выбрал строгий или нестрогий знак.
\item[$\square$] Я записал границы промежутка слева направо.
\item[$\square$] Я использовал круглые скобки рядом с $\pm\infty$.
\end{itemize}
\clearpage

\section*{4. Алгоритм: от рисунка к записи промежутка}
\headingrule
\vspace{2mm}
\begin{center}
\begin{tikzpicture}[node distance=10mm and 18mm,font=\small]
\node[fcstart] (start) {Начало};
\node[fcact,below=of start] (bounds) {Найдите выделенную часть оси. Запишите её левую и правую границы.};
\node[fcdecision,below=13mm of bounds] (inf) {Граница --- $\pm\infty$?};
\node[fcact,right=23mm of inf,minimum width=30mm,text width=29mm] (roundinf) {Круглая\\скобка};
\node[fcdecision,below=16mm of inf] (fill) {Точка на границе закрашена?};
\node[fcact,left=12mm of fill,minimum width=32mm,text width=31mm] (square) {Квадратная\\скобка};
\node[fcact,right=12mm of fill,minimum width=32mm,text width=31mm] (round) {Круглая\\скобка};
\node[fcact,below=19mm of fill,minimum width=66mm,text width=64mm] (write) {Запишите границу с выбранной скобкой. Повторите для другой границы.};
\node[fcstart,below=of write] (end) {Получили $x\in\text{промежуток}$};
\draw[fcline] (start)--(bounds);
\draw[fcline] (bounds)--(inf);
\draw[fcline] (inf)--node[above,font=\scriptsize]{да}(roundinf);
\draw[fcline] (inf)--node[right,font=\scriptsize]{нет}(fill);
\draw[fcline] (fill)--node[above,font=\scriptsize]{да}(square);
\draw[fcline] (fill)--node[above,font=\scriptsize]{нет}(round);
\draw[fcline] (square.south)|- (write.west);
\draw[fcline] (round.south)|- (write.east);
\draw[fcline] (roundinf.east)--++(0.4,0)|-([yshift=-2mm]write.east);
\draw[fcline] (write)--(end);
\end{tikzpicture}
\end{center}
\vspace{4mm}
\key{Пример алгоритма.} На оси закрашена точка $-1$, а штриховка идёт вправо. Левая граница $-1$ включена, правая --- $+\infty$. Значит,
\[
x\ge -1,\qquad x\in[-1;+\infty).
\]
\clearpage

\section*{5. Дополнительный сюжет урока: текстовая задача}
\headingrule
\key{Дано.} У Маши и Васи сначала было поровну денег. Маша потратила 6~рублей, Вася --- 14~рублей. После этого у Маши осталось в 3 раза больше денег, чем у Васи. Сколько рублей было у каждого сначала?

\begin{center}
\renewcommand{\arraystretch}{1.33}
\begin{tabularx}{.87\linewidth}{@{}Xcc@{}}
\toprule
& \textbf{До покупок} & \textbf{После покупок}\\
\midrule
Маша & $x$ & $x-6$\\
Вася & $x$ & $x-14$\\
\bottomrule
\end{tabularx}
\end{center}

\key{Как перевести фразу «у Маши в 3 раза больше».} После покупок у Маши \emph{большая} сумма, у Васи --- \emph{меньшая}; поэтому
\[
\underbrace{x-6}_{\text{Маша}}
=3\cdot\underbrace{(x-14)}_{\text{Вася}}.
\]
Неравенство сумм здесь превращаем в \emph{точное равенство с множителем $3$}. Умножить нужно сумму Васи --- она меньше.

\key{Решение.}
\[
\begin{aligned}
x-6&=3x-42,\\
x-3x&=6-42,\\
-2x&=-36,\\
x&=18.
\end{aligned}
\]
\key{Проверка.} У Маши осталось $18-6=12$ рублей, у Васи --- $18-14=4$ рубля. Действительно, $12=3\cdot4$. Изначально у каждого было по \key{18 рублей}.

\key{Типичная ошибка.} Равенство $x-6=x-14$ неверно: деньги после покупок уже \emph{не равны}. Соотношение «в 3 раза больше» необходимо отразить сразу при составлении уравнения.

\section*{6. Ошибки, которые важно предупреждать}
\headingrule
\begin{itemize}
\item \key{Смешение скобок.} У выколотой конечной точки круглая скобка; у закрашенной --- квадратная.
\item \key{Перепутанные знаки.} Штриховка справа означает $>$ или $\ge$, слева --- $<$ или $\le$.
\item \key{«Включили бесконечность».} Записи $[-\infty;5]$ и $[5;+\infty]$ неверны.
\item \key{Перечислены лишь целые числа.} Промежуток может содержать дробные числа, например $4{,}01$.
\item \key{Ошибочный множитель в задаче.} Фраза «у Маши в 3 раза больше» означает $\text{Маша}=3\cdot\text{Вася}$.
\end{itemize}
\clearpage

\section*{Тренировочный блок. Путь А}
\headingrule
\textit{10 заданий по возрастанию сложности. Для задач 1--8 используйте изученную связь между осью, неравенством и промежутком; в задаче 10 обязательно составьте уравнение.}

\begin{enumerate}[label=\textbf{\arabic*.},itemsep=1.0em]
\item Запишите промежуток $x\ge -1$ с помощью знака $\in$ и изобразите его на числовой оси.
\item Запишите промежуток $x<5$ с помощью знака $\in$ и изобразите его на числовой оси.
\item Запишите неравенство, задающее промежуток $x\in(-3;+\infty)$. Укажите, подходит ли число $-3$.
\item Запишите неравенство, задающее промежуток $x\in(-\infty;2]$. Укажите, подходит ли число $2$.
\item Запишите в виде промежутка двойное неравенство $-2\le x\le 2{,}5$. Назовите обе включённые границы.
\item Запишите в виде промежутка двойное неравенство $4\le x<12$. Какие из чисел $4$, $11{,}9$ и $12$ ему принадлежат?
\item Запишите двойным неравенством $x\in(4;12]$. Составьте \emph{два примера} подходящих чисел, одно из которых --- целое, другое --- дробное.
\item По рисунку запишите \emph{и} двойное неравенство, \emph{и} промежуток:
\begin{center}
\axes{\draw[accent,line width=2.4pt] (-2.2,0)--(2.5,0);\dotfull{-2.2}\dotopen{2.5}\node[below,font=\small] at (-2.2,-.1) {$-3$};\node[below,font=\small] at (2.5,-.1) {$1$};}
\end{center}
\item Для $x\in(-1;3]$ выберите \emph{все} подходящие значения из списка:
$-1$; $-0{,}5$; $0$; $3$; $3{,}1$. Объясните выбор каждого из двух граничных значений $-1$ и $3$.
\item У Алины и Димы изначально было одинаковое количество денег. Алина потратила 8~рублей, Дима --- 20~рублей. После покупок у Алины осталось в 3 раза больше денег, чем у Димы. Сколько рублей было у каждого до покупок? Составьте уравнение и выполните проверку.
\end{enumerate}
\clearpage

\section*{Ответы к тренировочному блоку}
\headingrule
\begin{center}
\renewcommand{\arraystretch}{1.7}
\begin{tabularx}{\linewidth}{@{}c X@{}}
\toprule
\textbf{№} & \textbf{Ответ}\\
\midrule
1 & $x\in[-1;+\infty)$; точка $-1$ закрашена, луч вправо.\\
2 & $x\in(-\infty;5)$; точка $5$ выколота, луч влево.\\
3 & $x>-3$; число $-3$ не подходит.\\
4 & $x\le 2$; число $2$ подходит.\\
5 & $x\in[-2;2{,}5]$; включены $-2$ и $2{,}5$.\\
6 & $x\in[4;12)$; подходят $4$ и $11{,}9$; $12$ не подходит.\\
7 & $4<x\le 12$; например, $5$ и $4{,}5$.\\
8 & $-3\le x<1$; $x\in[-3;1)$.\\
9 & Подходят $-0{,}5$; $0$; $3$. Число $-1$ исключено, $3$ включено.\\
10 & У каждого по $26$ рублей. Уравнение: $x-8=3(x-20)$. Проверка: $18=3\cdot 6$.\\
\bottomrule
\end{tabularx}
\end{center}
\vspace{1.2em}
\key{Самопроверка.} В задачах с промежутками обязательно сверьте тип каждой скобки с условием. В текстовой задаче проверьте, какая сумма после покупок больше и во сколько раз.

\end{document}